\documentclass[11pt]{article}
\usepackage[margin=1in]{geometry}
\usepackage{amsmath,amssymb,amsthm}
\usepackage[hidelinks]{hyperref}

\newtheorem{theorem}{Theorem}
\newcommand{\egcd}{\operatorname{egcd}}

\title{Cubic counterexamples for perfect powers and their divisor sums}
\author{Pioneer project}
\date{28 November 2026}

\begin{document}
\maketitle

\begin{abstract}
The comments to OEIS A128607 and A128608 ask whether there are further
terms with prime-factor exponent gcd greater than two. We give witnesses
on both sides of the divisor-sum map. The least new input-side witness is
$424462145606577000=751530^3$, whose divisor sum is $1292054400^2$.
On the output side, the divisor sum of $21531558370^2$ is
$10773399^3$. The witnesses have short multiplicative certificates.
\end{abstract}

For $u>1$ let $\egcd(u)$ be the gcd of the exponents in its prime
factorization. The sequence A128607 lists the perfect powers $u$ for which
$\sigma(u)$ is a perfect power. Its commentary asks if any entries besides
$81$ and $343$ have $\egcd(u)>2$. The paired sequence A128608 gives the
corresponding values $\sigma(u)$ and asks for a higher power besides the
value $1729^3$. We use these numerical exceptions as given in the two
entries~\cite{oeis607,oeis608}.

\begin{theorem}
There is an $N>343$ with $\egcd(N)=3$ and $\egcd(\sigma(N))=2$.
The least $N>343$ for which both $N$ and $\sigma(N)$ are perfect powers
and $\egcd(N)>2$ is $424462145606577000$.
There is also an $M\ne 1857437604$ with $\egcd(M)=2$ and
$\egcd(\sigma(M))=3$, namely $463608005840717056900$.
\end{theorem}

\section{Witnesses}

For distinct primes $p_i$ the identity
\[
 \sigma\left(\prod_i p_i^e\right)=
 \prod_i (1+p_i+\cdots+p_i^e)
\]
follows by writing each divisor as a product of prime powers. The two
certificates below give the factorization of every factor in the right
side. The primality of the displayed bases can be checked by trial division
up to their square roots.

For $N=(2\cdot 3\cdot 5\cdot 13\cdot 41\cdot 47)^3=751530^3$ use
\[
\begin{array}{r|r|l}
p & 1+p+p^2+p^3 & \text{Factorization}\\ \hline
2 & 15 & 3\cdot5\\
3 & 40 & 2^3\cdot5\\
5 & 156 & 2^2\cdot3\cdot13\\
13 & 2380 & 2^2\cdot5\cdot7\cdot17\\
41 & 70644 & 2^2\cdot3\cdot7\cdot29^2\\
47 & 106080 & 2^5\cdot3\cdot5\cdot13\cdot17
\end{array}
\]
The product is
$2^{14}3^45^47^2 13^2 17^2 29^2=1292054400^2$. The input has all
prime-factor exponents equal to $3$. Its divisor sum has exponent gcd
$2$. The same calculation also gives another cubic input witness:
multiplying $N$ by the coprime cube $343$ multiplies its divisor sum by
$400=20^2$.

For $M=(2\cdot5\cdot7\cdot29\cdot37\cdot439\cdot653)^2=21531558370^2$
use
\[
\begin{array}{r|r|l}
p & 1+p+p^2 & \text{Factorization}\\ \hline
2 & 7 & 7\\
5 & 31 & 31\\
7 & 57 & 3\cdot19\\
29 & 871 & 13\cdot67\\
37 & 1407 & 3\cdot7\cdot67\\
439 & 193161 & 3\cdot31^2\cdot67\\
653 & 427063 & 7\cdot13^2\cdot19^2
\end{array}
\]
Here the product is $(3\cdot7\cdot13\cdot19\cdot31\cdot67)^3=10773399^3$.
Both roots are squarefree, giving the asserted gcds.

The certificates can be found by modular arithmetic on exponent vectors.
For the first, reduce the exponents of $(p+1)(p^2+1)$ modulo two and
find a nonempty binary dependence among the vectors. For the second,
factor $p^2+p+1$ and look for a subset whose vectors sum to zero modulo
three. In the latter case splitting a finite prime list into two halves
allows matching subset-sum vectors to their additive inverses.

\section{Minimality on the input side}

If $\egcd(u)>2$, then $u=b^e$ with either $e=4$ or $e$ an odd
prime. In fact, any integer greater than two either has an odd prime
divisor or is divisible by four. For $2\leq u\leq N$, an exponent $e$
satisfies $2^e\leq N$, hence $e\leq58$. For each of the relevant exponents,
enumerate $2\leq b\leq\lfloor N^{1/e}\rfloor$. The required base
factorizations are obtained by a least-prime-factor sieve up to
$751530=\lfloor N^{1/3}\rfloor$. Multiplying the exact geometric
sums yields $\sigma(b^e)$. Testing that number for being a perfect power
gives the following complete output, grouped by exponent.
\[
\begin{array}{r|r|r|r}
e & b & b^e & \sigma(b^e)\\ \hline
3 & 7 & 343 & 400\\
3 & 751530 & 424462145606577000 & 1669404572559360000\\
4 & 3 & 81 & 121
\end{array}
\]
The computation is reproduced by the accompanying C verifier, using
GMP integer roots, exact division, and its deterministic perfect-power
test. Duplicate representations $u=b^e$ would be harmless because every
higher-exponent term must be encountered at least once.

\begin{thebibliography}{9}
\bibitem{oeis607}
OEIS Foundation Inc., \emph{A128607: Perfect powers whose divisor sums are
perfect powers}, \url{https://oeis.org/A128607}.
\bibitem{oeis608}
OEIS Foundation Inc., \emph{A128608: Divisor sums of the terms of A128607},
\url{https://oeis.org/A128608}.
\end{thebibliography}
\end{document}
